\section{Introduction}

Quantum resource estimates are shaped by compilation choices.  Recent systems make hardware assumptions configurable, optimize across full stacks, or predict compilation options from circuit features \cite{campbell2026resourceestimation,autoqureo2026,quetschlich2023options,quetschlich2025predictor}.  These developments improve the usefulness of estimates, but they leave a logically prior question: why should a particular numerical prediction be believed?  A formula can be correct because a theorem forces it, because a finite benchmark happened to cover the input, or because a structural regime model assigned the right label.  The same number can carry any of these warrants.

Verified optimizers establish transformation correctness, while exact synthesis and pattern matching can establish bounded optimality or locally exact rewrites \cite{hietala2021verified,li2022oracles,iten2022patterns,bravyi2022clifford}.  Algorithm selection and instance-space analysis instead ask which method works in which part of a feature space \cite{rice1976algorithm,smithmiles2009meta,smithmiles2023isa}.  Neither tradition licenses collapsing theorem authority, finite agreement, structural explanation and transfer into one accuracy score.  That collapse is especially risky when an interpretable formula survives extensive testing before one exact counterexample defeats it.

We study this separation in an exact shared-tag Pauli block-encoding model derived from Tag-and-Restore Encoding (TARE) \cite{schillo2026tare}.  The model admits unrestricted dynamic-programming optima on finite instances and an all-size support-two normalization theorem under a frozen structural objective.  This combination makes it possible to distinguish four questions cleanly.  Is the predicted cost exact?  Does the implementation realize the proved restricted family?  Does a compact structural vocabulary explain which minimizing configuration is used?  Does the same representation transfer to a materially different compiler family?

The evidence forms one refutation--repair--transfer argument.  A compact three-family forecaster first succeeds prospectively, then fails exactly on a fresh three-qubit instance.  Replacing that formula with the theorem-backed support-two optimum repairs the cost layer without deleting the counterexample.  Hostile tests then refute two successive explanation vocabularies while leaving the cost theorem intact.  A complete per-block-permutation analysis closes the previously open commuting-support-two sector, but it does not convert two failed alternative proof chains into successes.  Finally, weighted Clifford state preparation shows that even exact finite feature determination need not transfer to the next system size.

The resulting contribution is a layered certificate discipline supported by model-specific theorems and adverse tests.  It is not a claim that regime mapping, algorithm selection, Pauli compilation, verified optimization or quantum resource estimation begins here.  Nor are the compiler models treated as a statistical sample from which prevalence can be inferred.  Their role is discriminating: each exposes a concrete inference that would be valid only if different evidence types were mistakenly conflated.
